\section*{АННОТАЦИЯ}
\addcontentsline{toc}{section}{Аннотация}

Данная выпускная квалификационная работа посвящена разработке интеллектуальной системы прогнозирования страховых рисков и динамического расчета страховых тарифов в автостраховании на основе методов машинного обучения.

В ходе работы был проведен анализ традиционных методов, нацеленных на оценку рисков и обоснована необходимость перехода к вероятностной модели прогнозирования. В качестве основополагающего алгоритма был выбран CatBoost, который обеспечивает высокую точность в работе с категориальными признаками и гетерогенными данными. Разработана математическая модель бинарной классификации и формула динамического расчета коэффициента бонус-малус на основе предсказанной вероятности ДТП. Реализованы программные модули предобработки данных, обучения модели и расчета страхового тарифа с веб-интерфейсом, а тестирование модели метриками качества классификации подтвердило её предсказательную способность и чувствительность к ключевым факторам риска.

Разработанное решение может быть применено в сфере автострахования для внедрения персонализированного ценообразования, которое повысит точность оценки рисков и оптимизирует тарифную политику данной области.